\scanpage{37}\mspage{36}
% L1–2
plus petit possible est $Y=\overline X$, d'où \deleted{$Z=\overline X\unsure{,|\wedge}$}
\[
\overline X^{(\mathcal M_A)}=\overline X^{\mathcal M}\wedge A
\]
% L3–6
Soient $X,Y\in \mathcal M_{\deleted{X}A}$, $\{X,Y\}$ modéré) alors
\[
\begin{aligned}
\overline{X\deleted{\unsure{\wedge}}\vee Y}^{\mathcal M_A}
   &=\underbrace{\overline{X\vee Y}^{\mathcal M}\wedge A}\\
   &\hspace{-.4em}(\overline X\vee\overline Y)\wedge A
       =\underbrace{(\overline X\wedge A)}\vee(\overline Y\wedge A)\\
   &\hspace{12em}\overline X^A\vee\overline Y^A
\end{aligned}
\]

% L7–11
7) $X,Y\in \mathcal M_{\inserted{A}}^{\deleted{A}}$, $\{X,Y\}$ modéré) \deleted{\unsure{alors|donc}} \deleted{donc $\unsure{\{\overline X^{\mathcal M},\overline Y^{\mathcal M}\}}$ \unsure{aussi}} $\{X,Y,A\}$ mod. donc $\{\overline X,Y,A\}$ mod donc $\{\overline X\wedge A=\overline X^A,Y\{$ modéré.

\uline{OK}

% L12–15
\uline{Composantes}\qquad On dit qu'un $\mathcal M$ est connexe si \deleted{$\forall\,X$ fermé et ouvert,} $\forall\,X\in \mathcal M$, $X$ fermé et $X$ ouvert
\[
\Longrightarrow\quad X=\varnothing\quad\text{ou}\quad X=\underline{1}
\]

% L16–21
\deleted{\unsure{N|S}}\uline{NB}\qquad Soit $X$ tel que $X$ fermé et ouvert, On s'attend qu'il soit mod. avec \deleted{\uline{\unsure{$\forall A$|$A$}}} $Y$. S'il en est ainsi, Alors les $X$ en question engendrent une algèbre de Boole. \deleted{Si elle est finie, \unsure{les|on} \unsure{points|peut}} \deleted{\unsure{du spectre}} \inserted{les éléments de} son spectre jouent le rôle de comp. conn. Si l'alg.
